Understanding when clinical events occur is fundamental to modeling patient trajectories, enabling tasks such as process mining, outcome forecasting, and causal inference in medicine \citep{jensen2012mining, noren2010temporal}.
However, most clinical narratives lack explicit temporal structure. While electronic health record (EHR) datasets like MIMIC-III and MIMIC-IV include timestamped structured data and clinical notes \citep{johnson2016mimic, johnson2023mimic}, they capture only a subset of events (e.g., vital signs, lab results) and still require sophisticated natural language processing to extract temporal information from unstructured text \citep{viani2021temporal, cheng2023typed}.

In contrast, published case reports often provide rich narrative descriptions of patient timelines, but the timing of events is typically expressed in relative terms within the free text (e.g., “on day 3 of hospitalization”).
Efforts in temporal reasoning on clinical text have been hampered by the scarcity of large-scale annotated corpora.
Earlier initiatives, such as the 2012 i2b2 challenge and Clinical TempEval, focused on temporal relation extraction from clinical notes \citep{sun2013evaluating, bethard2016semeval, galvan2018investigating} but were constrained by small datasets (e.g., 310 documents) from single institutions, limiting model performance and generalizability.
Systems that rely on metadata timestamps (e.g., admission or discharge dates) fail to capture the fine-grained sequence of clinical events described in text, which is essential for nuanced understanding and prediction.
There remains a pressing need for large-scale, temporally annotated corpora to support the development and evaluation of models for clinical temporal reasoning.

In this work, we present the PubMed Open Access Textual Time Series (PMOA–TTS) corpus, a new large-scale dataset comprising 124,699 openly available clinical case reports, each annotated with a structured timeline of events.
Leveraging large language models (LLMs)—specifically Llama 3.3 \citep{grattafiori2024llama} and DeepSeek R1 \citep{guo2025deepseek}—we transform each case report into a sequence of (clinical event, time) tuples that capture key clinical events and their relative timing within the patient narrative.
To our knowledge, PMOA–TTS constitutes the largest publicly available collection of clinical narratives with explicit temporal event structure.
We demonstrate that these temporal annotations enable novel forms of analysis and significantly enhance downstream modeling of patient outcomes.

Building PMOA–TTS requires addressing several challenges: identifying relevant case reports at scale, extracting timeline events with high fidelity, and rigorously evaluating the quality of the generated annotations.
Rather than relying on naive metadata filters, our LLM-based pipeline accurately identifies single-patient case narratives and resolves ambiguous free-text time expressions to extract fine-grained timelines.
We further enrich each case with structured patient demographics and diagnostic labels using prompt-driven extraction.
Finally, we showcase the practical value of PMOA–TTS in a downstream survival analysis task, where models leveraging the extracted textual time series demonstrate strong performance in early outcome prediction.


\textbf{Contributions.} In summary, this paper makes the following contributions:
\begin{itemize}[leftmargin=*,itemsep=0pt,parsep=0pt,partopsep=0pt,topsep=0pt]
    \item We introduce the PMOA Textual Time Series (PMOA–TTS) corpus, comprising 124,699 single-patient case reports from PubMed Open Access, each annotated with structured clinical timelines extracted using large language models (LLMs).
    \item We develop a scalable LLM-based pipeline for data extraction and annotation, including automatic identification of case reports, prompting strategies for timeline, demographic, and diagnosis extraction, and a manual evaluation framework to assess annotation quality.
    \item We quantitatively evaluate the accuracy of the extracted timelines, demonstrating high temporal ordering performance (concordance index $>0.9$) and introducing a novel Area Under the Log-Time Curve (AULTC) metric for analyzing timing precision and errors.
    \item We examine demographic distributions, frequently occurring diagnoses, and diagnosis co-occurrence patterns in the corpus, highlighting that it has wide clinical and demographic coverage.
    \item We show how to use the textual time series in downstream modeling by framing a survival prediction task, where incorporating extracted temporal event sequences improves predictive performance.
\end{itemize}

PMOA–TTS fills a critical gap by providing a large-scale, openly available resource for temporal reasoning in clinical narratives.
We hope this corpus will enable the development and benchmarking of new models for timeline extraction, temporal question answering, and longitudinal outcome prediction in biomedical NLP.

